\section{현장 활용방안}

C-BAT는 월 최대를 새로 쓸 위험이 큰 날을 당일 00:00에 알리는 데 쓴다. 기본요금은 월 최대를 갱신하는 소수의 날이 정하기 때문이다. 이 장은 요금을 정하는 날과 경보 규칙을 정의하고, 규칙을 과거 날짜에 적용해 실제 관측으로 채점한 결과와 경보일 점검 항목을 제시한다. 피크 예측의 합산 성능은 4장의 개발+9월 평가이고, 이 장의 경보 백테스트는 개발 구간 중 요금시간대가 있는 44일에 한정한다. 모형의 반사실 일정 이동 절감액, 경보 포착 사례, 실제 조치 후 절감액은 서로 구별한다.

\subsection{요금을 정하는 날과 경보 규칙}

기본요금은 매달 요금시간대의 최대수요와 래칫 하한 가운데 큰 값으로 정해진다. 요금시간대는 일요일·공휴일을 뺀 09--23시의 중간·최대부하 시간대다. 래칫 하한은 지난 1년 동안 래칫 적용월(12·1·2·7·8·9월)의 최대수요에서 온다. 따라서 청구 수요를 움직이는 날은 그 달의 최대를 새로 쓰는 날뿐이다.

개발 구간(2021년 7월 7일--8월 31일) 가운데 요금시간대 관측이 있는 44일에서 월 최대를 갱신한 날은 7월 19일 하루였다. 그날 요금시간대 최대는 222 kW로, 전날까지의 기준선 206 kW를 16 kW 넘었다. 기본요금 8,320원/kW·월을 적용하면 그 달 기준 133,120원이다. 7월은 래칫 적용월이므로 이 초과분은 이후 최대 11개월의 하한에도 영향을 줄 수 있다. 연 단위로는 133,120원 $\times$ 12 $\approx$ 160만 원이 상한이며, 이후 달의 실제 피크가 높으면 줄어든다. 실제로 8월의 요금시간대 최대는 207 kW였지만 청구 수요는 7월이 만든 래칫 222 kW였다.

경보 규칙은 백테스트 결과를 보기 전에 정해 기록하였다. 예측일 $d$의 00:00에 다음 절차를 실행한다.
\begin{enumerate}
\item 기준선 $B_d$를 구한다. $B_d$는 이번 달 1일부터 전날까지의 요금시간대 관측 최대와, 이번 달 이전 래칫 적용월의 요금시간대 관측 최대 가운데 큰 값이다. $d$일 이전 관측만 쓴다.
\item C-BAT 사후 경로 $s=1,\dots,S$마다 요금시간대 슬롯의 최댓값 $M_d^{(s)}$를 구한다.
\item 갱신 확률 $P_d$와 기대 초과 기본요금 $E_d$를 계산한다.
\begin{equation}
P_d = \frac{1}{S}\sum_{s=1}^{S} \mathbf{1}\bigl[M_d^{(s)} > B_d\bigr], \qquad
E_d = 8{,}320 \times \frac{1}{S}\sum_{s=1}^{S} \max\bigl(0,\ M_d^{(s)} - B_d\bigr).
\end{equation}
시드 3개(0·1·2)의 평균을 쓴다.
\item $E_d$가 점검 비용 3만 원(담당자 1--2시간 가정) 이상이면 경보한다. 1만 원, 10만 원, $P_d \ge 0.5$는 민감도로 본다.
\end{enumerate}

위험은 Platt 보정을 거치지 않은 경로 비율을 그대로 쓴다. 7일 내부 보정은 보정 구간에 C90 초과일이 없으면 절편만 추정되어 모든 날의 위험을 같은 값으로 만든다. 개발 구간에서도 C90 판별력은 보정 전 AUC 0.926, 7일 Platt 보정 뒤 0.788로 보정 전이 높았다.

\subsection{백테스트와 경보일 조치}

주 규칙은 유일한 갱신일인 7월 19일을 잡았고, 44일 동안 10번 경보했다(표~\ref{tab:ch5-alarm}). 기준일을 그대로 쓰는 B0\_kind는 이날을 놓쳤다. 표의 금액은 갱신일의 피크를 기준선 아래로 완벽하게 막았다고 가정한 그 달 기준 상한이다.

\begin{table}[h]
\centering
\caption{월 최대 갱신 경보 백테스트(2021년 7월 7일--8월 31일 중 요금시간대가 있는 44일, 갱신일 1일).}
\label{tab:ch5-alarm}
\small
\begin{tabular}{lrrrr}
\toprule
규칙 & 갱신일 잡음 & 경보 수 & 오경보 & 금액 상한(원) \\
\midrule
$E_d \ge$ 1만 원 & 1/1 & 23 & 22 & 133,120 \\
$E_d \ge$ 3만 원(주 규칙) & 1/1 & 10 & 9 & 133,120 \\
$E_d \ge$ 10만 원 & 1/1 & 3 & 2 & 133,120 \\
$P_d \ge 0.5$ & 1/1 & 3 & 2 & 133,120 \\
C-BAT 점예측 $> B_d$ & 1/1 & 3 & 2 & 133,120 \\
B0\_kind 기준일 $> B_d$ & 0/1 & 0 & 0 & 0 \\
매일 경보 & 1/1 & 44 & 43 & 133,120 \\
\bottomrule
\end{tabular}
\end{table}

C-BAT는 7월 19일에 갱신 확률 0.58, 기대 초과 비용 약 12.3만 원(122,919원)을 냈다. 경로 최댓값의 중앙값은 211 kW였다. B0\_kind의 기준일 곡선은 이날 요금시간대 최대를 195 kW로 두어 기준선에 못 미쳤다. 점검 비용을 10만 원으로 올리면 경보는 3번으로 줄고 갱신일은 여전히 잡힌다.

주 규칙의 오경보 9건은 7월 초(7월 7--16일)에 몰렸다. 이 기간에는 기준선 206 kW가 평소 가동일 피크와 가까웠고, 여기에 C-BAT가 평범한 가동일의 피크를 높게 예측하는 경향이 겹쳤다. 갱신일이 1일뿐이므로 이 결과는 성능 추정이 아니라 사례다.

\subsubsection{금액 상한과 점검 비용}

완벽한 조치를 가정한 포착 금액 상한은 갱신일의 초과분을 합한 값이며 실제 절감액은 아니다. 주 규칙의 경보일 집합을 $\mathcal A$, 실제 갱신일 집합을 $\mathcal E$라 하면 그 달 기준 비교는 다음과 같다.
\begin{align}
U_{\mathrm{month}} &= 8{,}320\sum_{d\in\mathcal A\cap\mathcal E}
(M_d^{\mathrm{band}}-B_d)^+=133{,}120\ \text{원},\\
K_{\mathrm{inspect}} &= c|\mathcal A|=30{,}000\times10=300{,}000\ \text{원},\\
U_{\mathrm{month}}-K_{\mathrm{inspect}} &= -166{,}880\ \text{원}.
\label{eq:ch5-month-net-cap}
\end{align}
이는 주 규칙에서 가정한 점검비를 모든 경보일에 실제 지출하고, 포착한 갱신일의 초과분은 전부 막는 경우의 계산이다. 래칫의 이후 달 영향과 조치 비용은 포함하지 않는다. 따라서 이번 당월 계산만으로 순편익이 양수라고 주장할 수 없다. 앞서 제시한 연 약 160만 원도 이후 피크에 따라 달라지는 래칫 상한이며 관측된 연간 절감액이 아니다. 현장에서는 실제 점검 시간, 조치 비용과 후속 월 청구 수요를 기록해 순편익을 평가해야 한다.

경보일 점검 항목은 관측 조건을 바탕으로 제안하지만 저감 효과는 아직 검증하지 않았다. 최종 적용은 설비·작업 순서·납기를 아는 담당자가 정한다.
\begin{itemize}
\item 초반 부하 점검. 유사일 비교에서 초반 두 시간 생산 비중이 10\%p 높은 날은 요금시간대 피크가 9.3 kW 높았다. 담당자가 실제 설비 가동 기록으로 동시 기동 여부를 확인하고, 작업 제약이 허용하면 기동 순서를 분산하는 조치를 시험한다.
\item 08시·11시 부하 집중 점검. 개발 구간 고피크일 13일 가운데 10일이 08시나 11시에 피크를 기록했다. 이 집계만으로 동시 기동이 원인이라고 단정할 수는 없다. 11시는 여름철 최대부하 시간대로 기본요금 산정에 들어간다.
\item 기준선 근접 시 비필수 부하 지연. 당일 요금시간대 부하가 $B_d$에 가까워지면 비필수 부하를 뒤로 미룬다.
\end{itemize}

조치의 방향은 모형과 무관한 유사일 비교에서 왔다. 2021년 9월 이전 가동일 가운데 가동유형과 날짜유형이 같고 일 생산량 차이가 10\% 이내인 날로 쌍을 만들고 시작 시각이나 초반 생산 비중만 다른 쌍의 실측 피크를 비교했다. 시작이 늦은 날의 쌍 206개에서는 늦게 시작한 날의 요금시간대 피크가 더 높았다. 피크 차이의 합을 시작 지연의 합으로 나누면 1시간당 7.4 kW였다(95\% 구간 $[3.4,\ 12.0]$). 그러나 쌍별 피크 차이의 중앙값은 2.0 kW로, 큰 차이는 일부 쌍에 몰려 있다. 초반 두 시간 비중이 다른 쌍은 25개뿐이었고 10\%p당 9.3 kW(95\% 구간 $[0.6,\ 20.2]$)였다. 한 날이 여러 쌍에 들어가므로 구간은 실제보다 좁을 수 있고, 시작 시각은 수주량이나 설비 상태와 함께 움직일 수 있다. 따라서 이 값은 관측된 연관이며 조치의 효과 크기로 쓰지 않는다.

\subsection{모형 반사실을 쓰지 않는 이유와 현장 검증}

C-BAT로 생산계획만 바꾼 반사실 곡선을 계산하면 일정 이동의 절감액을 낼 수 있지만, 그 금액은 관측으로 뒷받침되지 않았다. 앞의 유사일 231쌍에서 두 날의 실측 차이와, 두 날의 계획을 서로 바꿔 넣은 모형 차이를 비교했다. 전력량요금 차이의 상관은 $-0.01$이었고 평균 방향도 반대였다. 관측에서는 시작이 늦은 날이 하루 7,359원 비쌌지만 모형은 464원 싸다고 예측했다. 요금시간대 피크에서는 모형의 반응이 관측 차이의 약 1/10이었다(관측 $+11.0$ kW, 모형 $+0.2$ kW). 날짜 단위 cluster bootstrap 95\% 구간은 0을 포함했다. 관측 쌍은 계획 외의 조건도 다를 수 있어 인과 기준값은 아니다. 그래서 모형이 틀렸다고 결론짓지는 않으며 모형 절감액을 관측으로 뒷받침할 수 없다는 데서 멈춘다.

실제 기본요금 규칙을 적용하면 기존의 일정 이동 권고는 기본요금을 바꾸지 못한다. 7월 청구 수요는 7월 19일 하루가 정했는데, 권고일에 이날은 없었다. 8월은 7월이 만든 래칫 222 kW가 청구 수요를 정해, 8월 피크를 얼마나 낮추든 기본요금이 같다. 권고의 기본요금 변화는 7·8월 모두 0원이었고, 모형의 피크 감소를 10배로 키워도 0원이었다. 남는 전력량요금 절감은 위 검증에서 지지되지 않았다. 이 때문에 절감 주장은 관측된 요금 구조, 즉 갱신일 하루가 만든 금액 상한까지로 제한한다.

조치의 효과는 현장 운영으로 검증한다. 작업 제약이 허용하는 경보일을 사전에 정한 무작위 배정으로 조치일과 비조치일에 나누고, 가동유형과 계획 생산량을 함께 기록해 두 집단의 요금시간대 피크를 비교한다. 갱신일이 드물기 때문에 비교 지표는 갱신 여부보다 경보일의 요금시간대 피크를 우선한다. 실제 점검비와 조치 비용을 포함한 순편익도 기록한다. 관측 연관성과 효과를 구별하려면 임의로 조치하기 쉬운 날만 고르지 않고 배정·미이행 사유를 남겨야 한다.

도입은 다음 순서로 진행한다.
\begin{enumerate}
\item 매일 00:00에 당일 생산계획으로 C-BAT를 실행해 $B_d$, $P_d$, $E_d$, 피크 분포와 가장 가능성이 높은 피크 시각을 담은 일일 보고를 만든다.
\item 담당자가 경보일 보고를 확인하고 설비·작업 순서·납기를 고려해 위 조치 가운데 적용할 것을 정한다. 적용 여부와 내용을 기록한다.
\item 매월 말에 경보 수, 오경보, 갱신일 포착 여부, 조치일과 비조치일의 요금시간대 피크를 검토하고 점검 비용 기준을 조정한다.
\end{enumerate}
